% =====================================================================
\newpage
\section*{Заключение}
\phantomsection
\addcontentsline{toc}{section}{Заключение}

В ходе выполнения лабораторной работы был реализован двумерный клеточный автомат с окрестностью фон Неймана, тороидальными граничными условиями и правилом перехода 20\,874\,436. Программа позволяет задавать ширину поля и количество итераций, выбирать начальные условия (случайные с заданной плотностью и \texttt{seed}, ручной ввод, одна клетка в центре или пустое поле), а также автоматически обнаруживает цикл и вычисляет его предпериод и период.

Анализ показал, что рождение клетки возможно только при живом западном соседе, поэтому у динамики есть выделенное направление --- структуры смещаются вправо, а одиночная клетка и диагональ «/» из трех клеток являются глайдерами с периодом $n$ на торе. Конфигурации, занимающие одну строку, эволюционируют по одномерному элементарному правилу 120. Строки без живых клеток остаются пустыми, поэтому паттерны растут только по горизонтали. Случайные конфигурации остаются неупорядоченными. Длины циклов быстро растут с размером поля: при $n \le 18$ циклы находятся и их периоды достигают десятков тысяч шагов, а при $n \ge 20$ за $10^5$ шагов цикл не обнаружен ни в одном запуске. По классификации Вольфрама автомат отнесен к классу III (хаотическое поведение).

\subsection*{Достоинства реализации}

\begin{itemize}
	\item Случайные начальные условия воспроизводимы благодаря \texttt{seed}, который выводится, даже если был сгенерирован автоматически.
\end{itemize}

\subsection*{Недостатки реализации}

\begin{itemize}
	\item Номер правила и порядок клеток в индексе заданы в коде, изменить их при работе программы нельзя.
	\item Поле выводится символами в консоль, поэтому большие поля ($n > 60$) неудобно анализировать.
\end{itemize}

\subsection*{Масштабирование}

\begin{itemize}
	% \item Поиск цикла алгоритмом Брента или Флойда, хранящим лишь несколько конфигураций вместо всех, что позволит исследовать периоды на больших полях.
	% \item Ввод номера правила и порядка клеток в индексе пользователем, поддержка прямоугольных полей и окрестности Мура.
	\item Графический интерфейс.
	\item Параллельное вычисление шага для больших полей.
\end{itemize}

